\subsection{Experimental Setup}
\label{sec:exp_setup}

\textbf{Benchmarks.}
The framework is evaluated on two temporally grounded VideoQA benchmarks, \textbf{NExT-GQA} \cite{xiao2024can} and \textbf{ReXTime} \cite{chen2024rextime}. NExT-GQA emphasizes visually grounded answering, requiring the model to predict the correct answer while localizing the supporting temporal evidence \cite{xiao2024can}. ReXTime focuses on reasoning across time, where the evidence required for answering may lie in a temporal region different from that implied by the question context \cite{chen2024rextime}.

\textbf{Evaluation metrics.}
For NExT-GQA, following prior work \cite{xiao2024can,liu2025videomind}, the reported metrics include both IoP-based and IoU-based grounding measures, namely mIoP, IoP@0.3, IoP@0.5, mIoU, IoU@0.3, and IoU@0.5, together with Acc@GQA (IoP$>$0.5) and Acc@QA. For ReXTime, following the benchmark protocol \cite{chen2024rextime}, the reported metrics are R@0.3, R@0.5, mIoU, Acc, and Acc@IoU=0.5.

\textbf{Baselines.}
On NExT-GQA, comparisons are made against representative grounded or temporally aware VideoQA methods, including FrozenBiLMNG+ \cite{yang2022zero}, VTimeLLM \cite{huang2024vtimellm}, SeViLA \cite{yu2023sevila}, LangRepo \cite{kahatapitiya2025language}, VideoStreaming \cite{qian2024streaming}, Grounded-VideoLLM \cite{wang2025grounded_videollm}, VideoExpert \cite{zhao2025videoexpert}, VideoChat-TPO \cite{yan2025tpo}, LLoVi \cite{zhang2024simple}, HawkEye \cite{wang2024hawkeye}, DeVi \cite{qin2025question_answering_dense_video_events}, and VideoMind \cite{liu2025videomind}. On ReXTime, comparisons are made against VTimeLLM \cite{huang2024vtimellm}, TimeChat \cite{ren2024timechat}, LITA \cite{huang2024lita}, and VideoMind \cite{liu2025videomind}. Unless otherwise noted, all baseline results are taken from the corresponding original papers.

\textbf{Implementation details.}
EviR contains four modules: a Refiner, a Retriever, a Validator, and a Generator. Gemini 2.5 Flash is used as the Refiner without additional training. The Retriever is initialized from Qwen3-VL-8B and trained by following the released TimeLens-8B recipe. Representative settings include one epoch of SFT with learning rate $1\times10^{-5}$, GRPO with learning rate $1\times10^{-6}$ and eight generations per prompt, and TimeLens video sampling at 2 FPS with up to 448 frames. The Validator is built on Qwen2-VL-7B and trained with LoRA-based supervised fine-tuning following the verifier design of VideoMind \cite{liu2025videomind}. Representative verifier settings follow VideoMind, including LoRA rank 64 with $\alpha=64$, one training epoch, and learning rate $5\times10^{-5}$. The Generator uses Qwen2-VL \cite{wang2024qwen2} without further fine-tuning. Table~\ref{tab:module_training} summarizes the trainable modules.



\begin{table}[!htbp]
\caption{Training configuration of the trainable modules in EviR.}
\label{tab:module_training}
\centering
\scriptsize
\setlength{\tabcolsep}{4pt}
\begin{tabular}{llll}
\hline
Module & Backbone & Training data & Training strategy \\
\hline
Retriever & Qwen3-VL-8B & TimeLens-100K & SFT $\rightarrow$ Filter $\rightarrow$ GRPO \\
Validator & Qwen2-VL-7B & DiDeMo, TACoS, QVHighlights & LoRA-based SFT \\
\hline
\end{tabular}
\end{table}
\subsection{Main Results}
\label{sec:main_results}

\begin{table}[t]
\caption{Results on NExT-GQA. PT denotes trained/pretrained models, and TF denotes training-free models. Best results are in bold.}
\label{tab:nextgqa_results}
\centering
\scriptsize
\setlength{\tabcolsep}{3pt}
\begin{tabular}{lccccccccc}
\hline
Method & Paradigm & mIoP & IoP@0.3 & IoP@0.5 & mIoU & IoU@0.3 & IoU@0.5 & Acc@GQA & Acc@QA \\
\hline
FrozenBiLMNG+ \cite{yang2022zero} & PT & 24.2 & 28.5 & 23.7 & 9.6 & 13.5 & 6.1 & 17.5 & 70.8 \\
VTimeLLM \cite{huang2024vtimellm} & PT & 27.9 & 30.3 & 23.8 & 18.3 & 27.7 & 14.1 & 12.7 & -- \\
SeViLA \cite{yu2023sevila} & PT & 29.5 & 34.7 & 22.9 & 21.7 & 29.2 & 13.8 & 16.6 & 68.1 \\
LangRepo \cite{kahatapitiya2025language} & PT & 31.3 & -- & 28.7 & 18.5 & -- & 12.2 & 17.1 & -- \\
VideoStreaming \cite{qian2024streaming} & PT & 32.2 & -- & 31.0 & 19.3 & -- & 13.3 & 17.8 & -- \\
Grounded-VideoLLM \cite{wang2025grounded_videollm} & PT & 34.5 & 42.6 & 34.4 & 21.1 & 30.2 & 18.0 & 26.7 & -- \\
VideoExpert \cite{zhao2025videoexpert} & PT & 34.6 & 45.3 & 29.3 & 27.9 & 41.0 & 22.4 & 21.6 & 71.1 \\
VideoChat-TPO \cite{yan2025tpo} & PT & 35.6 & 47.5 & 32.8 & 27.7 & 41.2 & 23.4 & 25.5 & -- \\
LLoVi \cite{zhang2024simple} & PT & 37.3 & -- & 36.9 & 20.0 & -- & 15.3 & 24.3 & -- \\
HawkEye \cite{wang2024hawkeye} & PT & -- & -- & -- & 25.7 & 37.0 & 19.5 & -- & -- \\
DeVi \cite{qin2025question_answering_dense_video_events} & TF & 39.7 & -- & 38.9 & 23.6 & -- & 19.5 & 28.9 & 73.1 \\
VideoMind \cite{liu2025videomind} & PT & 39.0 & \textbf{56.0} & 35.3 & 31.4 & \textbf{50.2} & 25.8 & 28.2 & 76.9 \\
\hline
\textbf{Ours} & PT & \textbf{41.76} & 54.13 & \textbf{39.02} & \textbf{34.86} & 49.88 & \textbf{31.57} & \textbf{31.05} & \textbf{77.0} \\
\hline
\end{tabular}
\end{table}



Table~\ref{tab:nextgqa_results} shows that EviR achieves the strongest overall grounding-aware performance on NExT-GQA, with the largest gains appearing on stricter grounding metrics. Compared with VideoMind~\cite{liu2025videomind}, EviR improves IoU@0.5 from 25.8 to 31.57 and Acc@GQA from 28.2 to 31.05, while also achieving the best Acc@QA. These results are consistent with the design of EviR: query reformulation produces retrieval queries that are better aligned with observable events, and evidence validation helps suppress near-miss temporal candidates before answer generation.


\begin{table}[t]
\caption{Results on ReXTime. FT indicates the fine-tuned setting reported by prior work. Best results are in bold.}
\label{tab:rextime_results}
\centering
\scriptsize
\setlength{\tabcolsep}{4pt}
\begin{tabular}{lcccccc}
\hline
Method & FT & R@0.3 & R@0.5 & mIoU & Acc & Acc@IoU=0.5 \\
\hline
VTimeLLM \cite{huang2024vtimellm} & \checkmark & 43.69 & 26.13 & 29.92 & 57.58 & 17.13 \\
TimeChat \cite{ren2024timechat} & \checkmark & 40.13 & 21.42 & 26.29 & 49.46 & 10.92 \\
VTimeLLM \cite{huang2024vtimellm} & -- & 28.84 & 17.41 & 20.14 & 36.16 & -- \\
TimeChat \cite{ren2024timechat} & -- & 14.42 & 7.61 & 11.65 & 40.04 & -- \\
LITA \cite{huang2024lita} & -- & 29.49 & 16.29 & 21.49 & 34.44 & -- \\
VideoMind \cite{liu2025videomind} & -- & 38.22 & 25.52 & 27.61 & \textbf{74.59} & 20.20 \\
\hline
\textbf{Ours} & -- & \textbf{46.91} & \textbf{35.94} & \textbf{36.99} & \textbf{74.59} & \textbf{27.25} \\
\hline
\end{tabular}
\end{table}

Table~\ref{tab:rextime_results} shows a similar trend on ReXTime. EviR achieves the strongest grounding performance while matching the best answer accuracy, improving R@0.5 from 25.52 to 35.94 and Acc@IoU=0.5 from 20.20 to 27.25 over VideoMind~\cite{liu2025videomind}. This improvement is particularly notable on ReXTime, where the supporting evidence is often temporally displaced from the most salient moment suggested by the question~\cite{chen2024rextime}.

Taken together, the results on NExT-GQA and ReXTime suggest that the main advantage of EviR is not simply higher answer accuracy, but better alignment between answer prediction and temporal evidence. On NExT-GQA, this appears as stronger performance on stricter grounding-aware metrics than on looser overlap thresholds. On ReXTime, the same design becomes particularly effective because the relevant evidence is often temporally displaced. These results support the central claim of this work that evidence-guided query reformulation and validation improve grounded video reasoning when answer text alone is a weak retrieval signal.

\subsection{Ablation Study}
\label{sec:ablation}

Table~\ref{tab:ablation_rextime_modules} presents an ablation study of the four components in EviR on ReXTime.

\begin{table}[!htbp]
\caption{Ablation results on ReXTime. Ref, Ret, Val, and Gen denote Refiner, Retriever, Validator, and Generator, respectively. The Generator-only variant does not produce temporal evidence spans, so grounding metrics are not reported.}
\label{tab:ablation_rextime_modules}
\centering
\scriptsize
\setlength{\tabcolsep}{4pt}
\begin{tabular}{ccccccccc}
\hline
Ref & Ret & Val & Gen & R@0.3 & R@0.5 & mIoU & Acc & Acc@IoU=0.5 \\
\hline
-- & -- & -- & \checkmark & -- & -- & -- & 68.00 & -- \\
-- & \checkmark & \checkmark & \checkmark & 33.01 & 25.19 & 27.17 & 73.24 & 19.22 \\
\checkmark & \checkmark & \checkmark & -- & 46.91 & 35.94 & 36.99 & 32.68 & 18.68 \\
\checkmark & \checkmark & \checkmark & \checkmark & 46.91 & 35.94 & 36.99 & 74.59 & 27.25 \\
\hline
\end{tabular}
\end{table}

Using only the Generator yields moderate answer accuracy (68.00) but does not produce temporal evidence, confirming that answer prediction without explicit grounding is insufficient for grounded VideoQA.

Introducing the Retriever and Validator enables temporal grounding and improves answer accuracy to 73.24. However, grounding performance remains limited (R@0.5 = 25.19, Acc@IoU=0.5 = 19.22), indicating that directly using answer candidates as retrieval queries leads to suboptimal evidence localization.

Incorporating the Refiner significantly improves grounding quality, increasing R@0.5 from 25.19 to 35.94 and mIoU from 27.17 to 36.99. This demonstrates that rewriting answer candidates into evidence-oriented queries is critical for aligning retrieval with visually grounded events.

Notably, removing the Generator while retaining strong grounding performance results in a sharp drop in answer accuracy (32.68). This reveals a key limitation of evidence ranking: selecting the highest-scoring evidence alone is insufficient for grounded multiple-choice VideoQA, as the top-ranked clip may not be sufficiently discriminative to distinguish between answer candidates that share similar entities or coarse visual context.

Finally, combining all components yields the best overall performance, achieving strong grounding metrics together with the highest answer accuracy. These results highlight the complementary roles of the four modules: the Refiner improves query--evidence alignment, the Retriever localizes candidate spans, the Validator ranks evidence reliability, and the Generator performs final answer reasoning based on the selected evidence.
